% !TEX root = tarasov-qasp-zh.tex

\chapter[阿基米德原理在宇宙飞船上成立吗？]{阿基米德原理在宇宙飞船上成立吗？\label{ch-17}}

\begin{dialogue}

\tchr 当宇宙飞船处于失重状态时，阿基米德原理（Archimedes' principle）还成立吗？

\stda 我看它不成立。阿基米德原理的实质在于：由于物体与液体的密度不同（当然是指体积相同的物体与液体），把二者提升到同样高度所需做的功不同。而在失重状态下，这两份功没有差别，因为提升物体所需的功与提升同样体积的液体所需的功都等于零。

如果从液体对浸在其中的物体的压强来考虑，也会得出同样的结论，因为浮力（buoyant force）来源于物体的下底面与上底面所受压强之差。在失重状态下，这个压强差消失了，浮力也就随之消失。我还可以补充一点：失重状态下“上”与“下”没有分别，因此根本无法指明物体的哪个底面是上面，哪个是下面。

因此，\hlblue{在失重状态下，浸没在液体中的物体不受浮力作用。}这就是说，阿基米德原理对这种状态不成立。
\pnote{学生甲这段“功的论证”，是上一章（\chap{ch-16}）从能量考虑推导阿基米德原理的回声：在液体中把体积为 $V$ 的物体提升高度 $H$，会有同体积的液体下沉填补，故所需能量为 $(\rho-\rho_{0})gVH$，比在真空中做这份功省去 $\rho_{0}gVH$——省下的正是浮力对应的功。失重时它同样化为零，这与学生甲的结论一致。}

\stdb 我不同意\textbf{学生甲}的最后结论。我确信，阿基米德原理对失重状态同样是成立的。让我们推理得更仔细些。我们不直接过渡到失重状态，而是先考察一部以某一加速度 $a$ 运动的电梯，这个加速度与重力加速度 $g$ 同方向。设 $a<g$。容易看出，在这种情形，浸没在液体中的物体受到的浮力是
\begin{equation}%
F= \rho_{0} (g-a)V
\label{eq-96}
\end{equation}
而物体排开的那部分液体的重量同样等于 $\rho_{0} (g-a)V$。于是，浮力仍然等于物体排开的液体的重量，也就是说，阿基米德原理成立。下一步，我们逐渐增大加速度 $a$，使它向 $g$ 逼近。根据式~\eqref{eq-96}，浮力将逐渐减小，但与此同时，与物体同体积的液体的重量也以完全相同的方式减小。换句话说，在加速度 $a$ 趋近于加速度 $g$ 的全部过程中，阿基米德原理始终成立。在极限 $a=g$ 时，便进入失重状态。这时浮力变为零，但物体排开的液体的重量也同样变为零。因此，没有什么妨碍我们断言：阿基米德原理对失重状态同样成立。我想用一个例子来说明我的论证。设一块软木漂浮在盛水的容器中。根据式~\eqref{eq-95}，软木块浸入水中的体积与整块软木的体积之比，等于软木的密度与水的密度之比，即
\begin{equation}%
\frac{v_{1}}{v} = \frac{\rho_{3}}{\rho_{0}}
\label{eq-97}
\end{equation}
再设这个容器放在电梯里，电梯开始以某一加速度 $a$ 下降。由于软木和水的密度都没有因此改变，式~\eqref{eq-97} 仍然成立。换句话说，电梯做加速运动时，软木块相对于水面的位置与没有加速度时一样。显然，在 $a=g$ 而达到失重状态的极限情形下，这个条件也不会改变。这样一来，由阿基米德原理确定的软木块相对于水面的位置，原来与电梯的加速度无关。在这种情形下，失重的“有”与“无”是无法区分的。
\pnote{式~\eqref{eq-96} 把初中浮力与高中超重失重接在了一起：阿基米德原理（浮力等于物体排开的液体所受的重力）属初中内容（人教版八年级下册第十章《浮力》）；电梯以加速度 $a$ 下降时，一切“重力效应”——重量、液体压强、浮力——都按同一因子 $(g-a)$ 缩放（对照人教版必修第一册第四章“超重和失重”一节）。这里的“重量”取\chap{ch-09}所说的第二种含义，即弹簧秤在电梯中测出的视重。学生乙的论证相当于把 $g$ 换成有效值 $(g-a)$，再令它趋于零。}

\tchr 应当说，你们两人的论证都有充分的根据。不过，我必须同意学生甲：阿基米德原理对失重状态不成立。

\stdb 那您就得驳倒我的证明。

\tchr 我正要试试。你的论证建立在两个要点上。第一点：在加速度 $a<g$ 时，物体在液体中上浮的方式完全遵从阿基米德原理；第二点：这一论断在极限情形 $a=g$、即达到失重状态时也应当成立。对第一点我没有异议，第二点我不同意。

\stdb 可是您总不能否认，失重状态下软木块也同样停在原处不动！而它的位置正是直接由阿基米德原理得出的。

\tchr 对，这没错。失重状态下，软木块确实同样停在原处不动。然而在这种状态下，它相对于液面的位置已经不再是阿基米德原理的结果了。用手指把它按进深处，它就悬停在你把它留下的那个深度上。反过来，只要差值 $(g-a)$ 哪怕再小而不为零，软木块就会升到水面上，按阿基米德原理所确定的位置漂浮。可见，\hlblue{失重状态与哪怕极其微弱的有重性之间，存在着根本的差别。}换句话说，在向失重状态过渡时，在“最后一瞬间”发生了一个突然的变化，也就是一次突变（jump），使整个局面发生了质的改变。
\pnote{这一幕中国学生有机会亲眼看到：2021 年 12 月 9 日“天宫课堂”第一课（神舟十三号乘组）的“浮力消失实验”里，乒乓球被放入水中后不浮出水面，悬在水中——与教师这里描述的软木块如出一辙。另外，真实航天器中的失重并非严格：残余加速度约为 $g$ 的百万分之一量级（即“微重力”），浮力并未绝对为零，但与表面张力相比已完全不重要，液体的行为由毛细现象支配——气泡不上升，浸没的物体也谈不上“上浮”。（空间站中的失重即\chap{ch-09}所说的自由落体——绕地球持续运行。）}

\stdb 可是这个突变是由什么引起的？它是从哪儿来的？要知道在我的推理里，加速度 $a$ 是平滑地趋近加速度 $g$ 的。

\tchr 这个突变与下述事实有关：在 $a=g$ 时出现了一种对称性（symmetry）——“上”与“下”的差别消失了，这一点学生甲刚才已经非常贴切地指出了。如果差值 $(g-a)$ 无穷小，但毕竟不等于零，问题中就存在一个由物理确定的“向上”的方向，浮力恰恰就沿这个方向作用。然而在 $a=g$ 时，这个方向消失了，所有方向在物理上都变得等价。我所说的突变就是指这个。对称性的破坏或出现，总是伴随着突变。
\pnote{“突变”指的不是浮力本身发生跳变——按式~\eqref{eq-96}，浮力随 $(g-a)$ 连续地趋于零。在 $(g-a)=0$ 处突变的，是“方向”在问题中的地位：只要 $(g-a)$ 不为零，哪怕无穷小，问题中就存在一个优越方向（“上方”），液面取向、压强分布与浮力方向都由它定出；恰在 $(g-a)=0$ 时所有方向才真正等价，“上浮”随之失去意义。“对称性的破坏或恢复总伴随质变”是相变理论与近代物理（如对称性自发破缺）中的普遍主题，这里是它最朴素的一个例子。}

\end{dialogue}

\cleardoublepage

\thispagestyle{empty}
\vspace*{-1.95cm}
\begin{figure*}[!ht]
\centering
\includegraphics[width=0.73\textwidth]{figs/sec/tarasov-qasp-sec-09.pdf}
\end{figure*}
\vspace*{-.4cm}
\begin{fullwidth}
{\textsf{\Large \hlred{从根本上说，现代物理学就是分子物理学。因此，哪怕只是通过理想气体这个最简单的例子，去了解物质分子动理论的基础，也是格外重要的。水的热膨胀的特殊性问题将单独讨论。气体定律将得到详细分析，并将应用于具体工程问题的求解。
}}}
\end{fullwidth}

% NEWTERMS: Archimedes' principle | 阿基米德原理 | 本章首现括注；术语本属第 16 章范围（并行翻译中），全书首现归属以成书为准，届时可撤一处括注
% NEWTERMS: buoyant force | 浮力 | 人教版初中八年级下册第十章《浮力》；首现归属同上待第 16 章定
% NEWTERMS: submerged / immersed | 浸没（浸在液体中） | 语境词
% NEWTERMS: upper / lower base | 上底面／下底面 | 圆柱体语境（承第 16 章推导）
% NEWTERMS: buoyed up / come up to the surface | 上浮 | 与"漂浮 floating"区分
% NEWTERMS: lift | 电梯 | 承第 4、9 章译名
% NEWTERMS: piece of cork | 软木块 | 密度符号 $\rho_3$、体积 $v$、$v_1$ 承第 16 章
% NEWTERMS: weight of the liquid displaced | 排开的液体的重量 | 人教版初中表述作"排开的液体所受的重力"
% NEWTERMS: weightness | 有重性 | 承第 9 章定名（原书自造词）
% NEWTERMS: weightlessness | 失重 | 承第 7（章尾格言）、9 章，本章沿用不重复括注
% NEWTERMS: symmetry | 对称性 | 本章核心概念，首现括注
% NEWTERMS: jump | 突变 | 教师用语，首现括注（jump）
% NEWTERMS: spaceship | 宇宙飞船 | 章题
% NEWTERMS: pressure | 压强 | 基础术语不括注；系统展开见后续流体静力学各章
% NEWTERMS: molecular physics | 分子物理学 | 章尾格言首现
% NEWTERMS: molecular-kinetic theory (of matter) | 分子动理论 | 章尾格言首现；人教版选择性必修第三册第一章章名
% NEWTERMS: ideal gas | 理想气体 | 章尾格言首现
% NEWTERMS: thermal expansion | 热膨胀 | 章尾格言首现（水的热膨胀特殊性）
% NEWTERMS: gas laws | 气体定律 | 章尾格言首现
